\documentclass[11pt]{article}
\usepackage{icmlko}
\begin{document}

\title{관문이 제일 못 맞히던 도메인을 두 배로 만드는 자료를 거부했다 --- 진단 다섯 노트가 찾아낸 것을 규칙이 물렸다}
\author{월드모델 실험실}
\date{노트 539 · 2026년 8월 4일}
\maketitle

\begin{abstract}
노트 538 은 ``시장팝업의 \texttt{target\_breadth} 값 단계를 쪼개라''고 처방했다.
원본 레코드 $647$ 건을 열어 보니 \textbf{쪼갤 게 없었다} --- 현행 축부터 라벨
$\rho = +0.088$($p = 0.19$)로 무신호였다. \textbf{대신 안 쓰던 필드가 있었다.}

\begin{center}
\begin{tabular}{lrrr}
\hline
후보 & 덮음 & 값 개수 & 라벨 \\
\hline
현행 target\_breadth(ip\_or\_collab) & $100\%$ & $2$ & $\rho = +0.088$ ($p = 0.19$) \\
체험요소 개수 & $60\%$ & $7$ & $\rho = -0.017$ \\
컨셉 설명 길이 & $100\%$ & $154$ & $\rho = +0.100$ ($p = 0.13$) \\
\textbf{category} & $100\%$ & $9$ & \textbf{크루스칼 $p < 0.0001$} \\
\hline
\end{tabular}
\end{center}

\texttt{category} 를 \texttt{grp}(덮음 $0\% \to 95\%$)에 채우니
\textbf{시장팝업 $\rho$ 가 $+0.1538 \to +0.3130$ 으로 두 배}가 되고
\textbf{판이 $+0.0048$($8/8$)} 올랐다 --- 이 사이클 최대 채택
(\texttt{DOMDROP} $+0.0043$)보다 크다. \textbf{그런데 관문이 기각했다.}
\end{abstract}

\begin{figure}[t]
\centering
\includegraphics[width=\textwidth]{figs/f1.pdf}
\caption{\textbf{가} --- 시장팝업 도메인. \textbf{나} --- 다섯 자.
\textbf{다} --- 판정과 그 근거.}
\end{figure}

\section{다섯 자}

\begin{center}
\begin{tabular}{lrrl}
\hline
자 & 차 & 부호 & 판정 \\
\hline
판 & $\mathbf{+0.0048}$ & $8/8$ & \textbf{$0$ 밖 양수} \\
날짜 통제 판 & $-0.0014$ & $1/8$ & \textbf{$0$ 밖 음수} \\
무리 안 판(공통 집합) & $-0.0025$ & $2/8$ & $0$ 을 문다 \\
KR 만화 & $-0.0017$ & $2/8$ & $0$ 을 문다 \\
비게임 앱 & $-0.0035$ & $2/8$ & $0$ 을 문다 \\
\hline
\end{tabular}
\end{center}

\textbf{노트 446 거부권}(판$\uparrow$ 날짜통제$\downarrow$ $\to$ 기각) 발동.
\textbf{노트 426 거울}도 집 밖 짝 둘이 모두 음수 부호라 미달이다.
\textbf{기각한다.}

\section{내 측정 결함 하나를 잡았다}

처음에 무리 안 판이 $-0.0155$($0/12$)로 크게 나빴다. 그런데 그 자는
\texttt{grp} 으로 쪼개 재는 것이라, \textbf{전 팔에서는 시장팝업이 grp 이 없어
제외되고 후 팔에서는 포함된다} --- \textbf{분모가 달랐다}(노트 $82 \cdot 90$
위반). 시장팝업은 어느 셈법에서도 가장 어려운 도메인이라 새로 들어오면 평균을
끌어내린다.

\begin{quote}
\textbf{같은 도메인 집합으로 다시 재니 $-0.0025$ 로 $0$ 을 문다.}
$-0.0155$ 의 $84\%$ 가 분모 결함이었다.
\end{quote}

\section{거부권의 근거는 사실이 아니다}

노트 446 의 거부권은 ``판이 오르는데 날짜 통제가 내려가면 그 이득은
\emph{날짜 구멍}''이라는 읽기에서 나왔다. 여기서는 그렇지 않다.

\begin{itemize}
\item \texttt{category} $\leftrightarrow$ 연도 크루스칼 $p = \mathbf{0.168}$
--- \textbf{연도를 안 실어 나른다.} 연도별 비중도 안정적이다(character
$23\!\sim\!31\%$ $\cdot$ fashion $15\!\sim\!40\%$).
\item \textbf{시장팝업은 애초에 날짜 통제 판 계산에 안 들어간다}(유보 연도
종류 부족). \textbf{시장팝업을 빼고 재도 $-0.0014$ 로 똑같다} --- 하락은
\emph{전부 이웃 열 도메인}이다.
\end{itemize}

\begin{quote}
\textbf{그래도 기각한다.} 사후에 거부권을 고쳐 통과시키면 그것이 바로
\emph{점수 고르기}다. \textbf{규칙과 증거가 어긋난 사실을 숨기지 않고 적는
것}이 이 노트가 하는 일이고, 거부권을 다시 볼지는 \emph{별도 사전 등록}으로
한다.
\end{quote}

\section{진단이 처음으로 자료를 찾아냈다}

이 결과가 나온 경로를 적어 둔다 --- 노트 $530$(시장팝업 $126$ 행이 $27$ 가지
설명으로 무너진다) $\to$ $534$(어느 셈법에서도 꼴찌이고 성과 폭 탓이 아니다)
$\to$ $537$(접전을 사는 것은 사람이 태깅한 속성이다) $\to$ $538$(그중
target\_breadth 인데 시장팝업에서 두 값뿐이다) $\to$ 여기.

\begin{quote}
\textbf{다섯 노트의 진단이 실제로 쓸 수 있는 자료 한 칸을 지목했고, 그것이
그 도메인을 두 배로 만들었다.} 이 사이클의 노트들이 ``성능이 아니라 진단''
이었는데, \textbf{진단이 성능이 되는 경로가 처음으로 닫혔다.}
\end{quote}

\section*{조항}

\begin{quote}
\textbf{자 하나에 도메인이 새로 참여하면 그것은 자가 바뀐 것이다.} 무리 안
판처럼 \emph{어떤 축이 있어야 계산되는} 자는, 그 축을 채우는 변경을 잴 때
\textbf{공통 도메인 집합}으로 다시 재야 한다.
\end{quote}

\section*{남은 것}

\begin{itemize}
\item \textbf{공유 \texttt{grp} 대신 시장팝업 전용 열로 넣으면 이웃이 안
다치나} --- 이웃 손해가 공유 열의 뜻이 바뀌어서라면 전용 열은 안 다칠 것이다.
그러면 다섯 자를 다 통과할 수 있다.
\item \textbf{노트 446 거부권을 다시 볼지} --- 별도 사전 등록. 지금 사례는
``판$\uparrow$ 날짜통제$\downarrow$ 인데 날짜 구멍이 아닌'' 첫 반례다.
\item \textbf{같은 무늬가 다른 도메인에도 있나} --- \texttt{grp} 덮음이 $0\%$
인 도메인이 시장팝업 말고 팝업 $\cdot$ 아이돌 $\cdot$ 만화 셋 더 있다.
\end{itemize}

\end{document}
